\section{Related work}
\label{sec:relatedwork}

% = PARAGRAPH 1 - Commonsense knowledge graphs

ConceptNet's goal was to collect the everyday knowledge computers need to understand what people were talking about~\cite{speer2012conceptnet, speer2017conceptnet}. According to the creators of ATOMIC, ConceptNet was too focused on general knowledge and did not capture wider inferential commonsense~\cite{sap2019atomic}. So, they developed ATOMIC to address the coverage limitations of ConceptNet by using a crowdsourcing approach. However, this crowdsourced data, having a noisy nature, led to ATOMIC having quality issues. To address this, \citet{hwang2021cometatomic} developed ATOMIC$^{2020}$, which was created using a more controlled data collection process and includes 23 relation types across three categories (social interactions, physical interactions, and event-centered interactions) to improve the quality and coverage of commonsense knowledge.

% = PARAGRAPH 2 - Neural knowledge models

The COMET framework~\cite{bosselut-etal-2019-comet} was developed to bridge the gap between static knowledge graphs and generative language models. It fine-tunes a pre-trained language model on a structured knowledge graph (ATOMIC) to generate commonsense inferences. COMET-BART, introduced by \citet{hwang2021cometatomic}, is an extension of the COMET framework that uses the BART architecture and is trained on the updated ATOMIC$^{2020}$ knowledge graph. The original paper also compares COMET-BART's performance to GPT-3~\cite{brown2020language}, a much larger language model. The comparison showed that fine-tuning on structured knowledge (COMET-BART) outperforms raw scale (GPT-3) in generating commonsense inferences.


% = PARAGRAPH 3 - Evaluation of text generation and your gap

The evaluation metric used in the original paper to assess COMET-BART's performance is BLEU~\cite{papineni2002bleu}, which is a standard tool for evaluating generative models. Other related automatic metrics include ROUGE~\cite{lin2004rouge}, METEOR~\cite{banerjee2005meteor}, CIDEr~\cite{Vedantam2015CVPR} and BERTScore~\cite{zhang2019bertscore}. Prior work evaluated COMET-BART's performance at the aggregate level across all relation types. In contrast, the gap in the literature that this report addresses is the lack of a finer-grained analysis of COMET-BART's performance across different relation types. This report applies the same BLEU metric on a per-relation basis to provide a more detailed understanding of COMET-BART's strengths and weaknesses in generating commonsense inferences.
